\section*{References}
\begingroup\small\setlength{\parindent}{-1.5em}\setlength{\leftskip}{1.5em}
Ackerman, T. A. (1989). Unidimensional IRT calibration of compensatory and noncompensatory multidimensional items. \textit{Applied Psychological Measurement, 13}(2), 113--127. \url{https://doi.org/10.1177/014662168901300201}

Akaike, H. (1987). Factor analysis and AIC. \textit{Psychometrika, 52}(3), 317--332. \url{https://doi.org/10.1007/BF02294359}

Ansley, T. N., \& Forsyth, R. A. (1985). An examination of the characteristics of unidimensional IRT parameter estimates derived from two-dimensional data. \textit{Applied Psychological Measurement, 9}(1), 37--48. \url{https://doi.org/10.1177/014662168500900104}

Baik, J., Ben Arous, G., \& Péché, S. (2005). Phase transition of the largest eigenvalue for nonnull complex sample covariance matrices. \textit{Annals of Probability, 33}(5), 1643--1697. \url{https://doi.org/10.1214/009117905000000233}

Chun, S. Y., \& Shapiro, A. (2009). Normal versus noncentral chi-square asymptotics of misspecified models. \textit{Multivariate Behavioral Research, 44}(6), 803--827. \url{https://doi.org/10.1080/00273170903352186}

Cudeck, R., \& Browne, M. W. (1992). Constructing a covariance matrix that yields a specified minimizer and a specified minimum discrepancy function value. \textit{Psychometrika, 57}(3), 357--369. \url{https://doi.org/10.1007/BF02295424}

Dijkstra, T. K. (1990). Some properties of estimated scale invariant covariance structures. \textit{Psychometrika, 55}(2), 327--336. \url{https://doi.org/10.1007/BF02295290}

Fukasaku, R., Hirose, K., Kabata, Y., \& Teramoto, K. (2025). Algebraic approach to maximum likelihood factor analysis. \textit{Psychometrika, 90}(5), 1732--1764. \url{https://doi.org/10.1017/psy.2025.10047}

Golubitsky, M., \& Schaeffer, D. G. (1985). \textit{Singularities and groups in bifurcation theory} (Vol.~1). Springer. \url{https://doi.org/10.1007/978-1-4612-5034-0}

Groskurth, K., Bhaktha, N., \& Lechner, C. M. (2025). The simulation-cum-ROC approach: A new approach to generate tailored cutoffs for fit indices through simulation and ROC analysis. \textit{Behavior Research Methods, 57}(5), Article 135. \url{https://doi.org/10.3758/s13428-025-02638-x}

Groskurth, K., Bluemke, M., \& Lechner, C. M. (2024). Why we need to abandon fixed cutoffs for goodness-of-fit indices: An extensive simulation and possible solutions. \textit{Behavior Research Methods, 56}(4), 3891--3914. \url{https://doi.org/10.3758/s13428-023-02193-3}

Hancock, G. R., \& Mueller, R. O. (2011). The reliability paradox in assessing structural relations within covariance structure models. \textit{Educational and Psychological Measurement, 71}(2), 306--324. \url{https://doi.org/10.1177/0013164410384856}

Heene, M., Hilbert, S., Draxler, C., Ziegler, M., \& Bühner, M. (2011). Masking misfit in confirmatory factor analysis by increasing unique variances. \textit{Psychological Methods, 16}(3), 319--336. \url{https://doi.org/10.1037/a0024917}

Horn, R. A., \& Johnson, C. R. (2013). \textit{Matrix analysis} (2nd ed.). Cambridge University Press.

Hu, L., \& Bentler, P. M. (1999). Cutoff criteria for fit indexes in covariance structure analysis. \textit{Structural Equation Modeling, 6}(1), 1--55. \url{https://doi.org/10.1080/10705519909540118}

Kenny, D. A., \& McCoach, D. B. (2003). Effect of the number of variables on measures of fit in structural equation modeling. \textit{Structural Equation Modeling, 10}(3), 333--351. \url{https://doi.org/10.1207/S15328007SEM1003_1}

Kroenke, K., Wu, J., Yu, Z., Bair, M. J., Kean, J., Stump, T., \& Monahan, P. O. (2016). Patient Health Questionnaire Anxiety and Depression Scale: Initial validation in three clinical trials. \textit{Psychosomatic Medicine, 78}(6), 716--727. \url{https://doi.org/10.1097/PSY.0000000000000322}

Lai, K., \& Green, S. B. (2016). The problem with having two watches: Assessment of fit when RMSEA and CFI disagree. \textit{Multivariate Behavioral Research, 51}(2--3), 220--239. \url{https://doi.org/10.1080/00273171.2015.1134306}

Lawley, D. N. (1956). Tests of significance for the latent roots of covariance and correlation matrices. \textit{Biometrika, 43}(1--2), 128--136. \url{https://doi.org/10.1093/biomet/43.1-2.128}

Lubbe, D. (2023). Advantages of using unweighted approximation error measures for model fit assessment. \textit{Psychometrika, 88}(2), 413--433. \url{https://doi.org/10.1007/s11336-023-09909-6}

McNeish, D., An, J., \& Hancock, G. R. (2018). The thorny relation between measurement quality and fit index cutoffs in latent variable models. \textit{Journal of Personality Assessment, 100}(1), 43--52. \url{https://doi.org/10.1080/00223891.2017.1281286}

McNeish, D., \& Wolf, M. G. (2023). Dynamic fit index cutoffs for one-factor models. \textit{Behavior Research Methods, 55}(3), 1157--1174. \url{https://doi.org/10.3758/s13428-022-01847-y}

Moshagen, M., \& Auerswald, M. (2018). On congruence and incongruence of measures of fit in structural equation modeling. \textit{Psychological Methods, 23}(2), 318--336. \url{https://doi.org/10.1037/met0000122}

Paul, D. (2007). Asymptotics of sample eigenstructure for a large dimensional spiked covariance model. \textit{Statistica Sinica, 17}(4), 1617--1642. \url{https://www3.stat.sinica.edu.tw/statistica/oldpdf/A17n418.pdf}

R Core Team. (2025). \textit{R: A language and environment for statistical computing} (Version 4.5.1) [Computer software]. R Foundation for Statistical Computing. \url{https://www.R-project.org/}

Reckase, M. D. (1979). Unifactor latent trait models applied to multifactor tests: Results and implications. \textit{Journal of Educational Statistics, 4}(3), 207--230. \url{https://doi.org/10.3102/10769986004003207}

Reise, S. P., Scheines, R., Widaman, K. F., \& Haviland, M. G. (2013). Multidimensionality and structural coefficient bias in structural equation modeling: A bifactor perspective. \textit{Educational and Psychological Measurement, 73}(1), 5--26. \url{https://doi.org/10.1177/0013164412449831}

Rodriguez, A., Reise, S. P., \& Haviland, M. G. (2016). Evaluating bifactor models: Calculating and interpreting statistical indices. \textit{Psychological Methods, 21}(2), 137--150. \url{https://doi.org/10.1037/met0000045}

Rönkkö, M., \& Cho, E. (2022). An updated guideline for assessing discriminant validity. \textit{Organizational Research Methods, 25}(1), 6--47. \url{https://doi.org/10.1177/1094428120968614}

Rosseel, Y. (2012). lavaan: An R package for structural equation modeling. \textit{Journal of Statistical Software, 48}(2), 1--36. \url{https://doi.org/10.18637/jss.v048.i02}

Savalei, V. (2021). Improving fit indices in structural equation modeling with categorical data. \textit{Multivariate Behavioral Research, 56}(3), 390--407. \url{https://doi.org/10.1080/00273171.2020.1717922}

Shi, D., Lee, T., \& Maydeu-Olivares, A. (2019). Understanding the model size effect on SEM fit indices. \textit{Educational and Psychological Measurement, 79}(2), 310--334. \url{https://doi.org/10.1177/0013164418783530}

Strachan, T., Cho, U. H., Ackerman, T., Chen, S., de la Torre, J., \& Ip, E. H. (2022). Evaluation of the linear composite conjecture for unidimensional IRT scale for multidimensional responses. \textit{Applied Psychological Measurement, 46}(5), 347--360. \url{https://doi.org/10.1177/01466216221084218}

Tipping, M. E., \& Bishop, C. M. (1999). Probabilistic principal component analysis. \textit{Journal of the Royal Statistical Society: Series B (Statistical Methodology), 61}(3), 611--622. \url{https://doi.org/10.1111/1467-9868.00196}

Watanabe, S. (1960). Information theoretical analysis of multivariate correlation. \textit{IBM Journal of Research and Development, 4}(1), 66--82. \url{https://doi.org/10.1147/rd.41.0066}

\endgroup
